Stringer, Pachitariu, Steinmetz, Carandini and Harris showed that a population code that maps a $d$-dimensional stimulus space differentiably to noise-free responses has an eigenspectrum decaying asymptotically faster than $n^{-(1+2/d)}$, reported exponents of 1.49, 1.65 and 3.43 for stimulus sets with $d = 8$, 4 and 1, fitted over ranks 11--500 (ranks 5--30 for 32 grating directions), and concluded that the code is about as high-dimensional as smoothness allows. That a finite window of ranks cannot measure an asymptotic exponent is elementary and partly anticipated; this note makes it quantitative for these stimulus sets. For codes built on bounded eigenfunctions, as on a torus, a known kernel-matrix bound implies that the finite-sample spectrum is continuous in the variance of the tail and blind to its rate of decay, so no estimator continuous in that spectrum can tell a differentiable code from a non-differentiable one. We place noise-free codes of fixed smoothness (Mat\'ern tuning) on the stimulus coordinates of all ten 8D and 4D stimulus sets. Codes exactly at the differentiability border give ranks 11--500 exponents from \WinBorderEightMin{} to \WinBorderEightMax{} at $d = 8$ and from \WinBorderFourMin{} to \WinBorderFourMax{} at $d = 4$, below the bound for short and above it for long tuning length scales, with infinitely many or 8,704 neurons and in whitened coordinates. Non-differentiable codes with Mat\'ern $\nu = 0.75$ also exceed the bound in the unwhitened coordinates at 2,800 stimuli, but not over ranks 101--500, not in whitened 8D coordinates and not always with fewer stimuli. For 32 directions a non-differentiable code reaches \PropWSmoothA{}. The tail exponent of the eigenmoment method of Pospisil and Pillow depends on the unresolved tail: spectra that share a broken power law with tail exponent 1.25 up to rank 500 give tail exponents from \MemeTailMinA{} to \MemeTailMaxA{}, and the eigenmoment misfit detects the difference rarely or not at all. The reported exponents are consistent with the bound but cannot show that the code satisfies it or lies close to it.
